\documentclass[11pt,a4paper]{article}
\usepackage[margin=2.4cm]{geometry}
\usepackage[T1]{fontenc}
\usepackage{booktabs}
\usepackage[font=small,labelformat=empty,skip=4pt]{caption}
\usepackage{array}
\usepackage{xcolor}
\usepackage{titlesec}
\usepackage{enumitem}
\usepackage{tcolorbox}
\tcbuselibrary{breakable}
\usepackage{graphicx}
\usepackage{amsmath}
\usepackage{amssymb}
\usepackage[hidelinks]{hyperref}

\definecolor{good}{HTML}{1B7F3B}
\definecolor{warn}{HTML}{B26B00}
\definecolor{bad}{HTML}{B3261E}
\definecolor{grey}{HTML}{666666}
\definecolor{rule}{HTML}{2B4C7E}

\newcommand{\code}[1]{\texttt{#1}}
\newcommand{\dn}{\textcolor{good}{\textbf{done}}}
\newcommand{\fl}{\textcolor{warn}{\textbf{in flight}}}
\newcommand{\ns}{\textcolor{grey}{not started}}
\newcommand{\done}{\textcolor{good}{\textbf{DONE}}}
\newcommand{\prog}{\textcolor{warn}{\textbf{IN PROGRESS}}}
\newcommand{\todo}{\textcolor{grey}{\textbf{NOT STARTED}}}
\newcommand{\wrong}{\textcolor{bad}{\textbf{FAILED, AND KEPT}}}

\titleformat{\section}{\large\bfseries\color{rule}}{\thesection.}{0.6em}{}
\titleformat{\subsection}{\normalsize\bfseries}{}{0em}{}
\setlist[itemize]{leftmargin=1.1em,itemsep=2pt,topsep=3pt}
\setlength{\parskip}{5pt}
\setlength{\parindent}{0pt}
\renewcommand{\arraystretch}{1.25}

% the phase box: what we set out to do
\newtcolorbox{phasebox}[1]{breakable,colback=white,colframe=rule!55,boxrule=0.7pt,
  arc=2pt,left=8pt,right=8pt,top=6pt,bottom=6pt,title={\bfseries #1},
  fonttitle=\color{white},colbacktitle=rule}

% the run box: what came back when we ran it
\newtcolorbox{runbox}[1]{breakable,colback=grey!4,colframe=grey!55,boxrule=0.7pt,
  arc=2pt,left=8pt,right=8pt,top=6pt,bottom=6pt,title={\bfseries #1},
  fonttitle=\color{white},colbacktitle=grey!75!black}

\title{\textbf{A zebrafish eyeball, moved by six muscles}\\[6pt]
\large Progress note --- written to be read, not decoded}
\author{}
\date{3 August 2026}

\begin{document}
\maketitle
\vspace{-2.2\baselineskip}

\begin{center}\small\color{grey}
\texttt{prototype/eye/}. Operational summary: \texttt{STATUS.md}. Plain English throughout.\\
Organised by phase, because phases are where we stop and you decide. Each phase is two boxes:
what we \emph{set out} to do, and what came back when we \emph{ran} it.
\end{center}

\vspace{4pt}
\hrule
\vspace{10pt}

%======================================================================
\section{What this is}

A deformable zebrafish eyeball in a bony orbit, rotated by six extraocular muscles that are
themselves contracting MLS-MPM bodies, expressed entirely as a Plexus\,2 \code{spec.yaml}.

Nothing applies a force ``to the eye''. A muscle shortens because an active stress acts along its
fibres; it cannot reel itself up because its origin end is anchored to bone; and it moves the globe
because its tendon end is embedded in the sclera and the two bodies share one background grid.
\textbf{The globe rotates because a muscle got shorter.} That is the whole mechanism.

\begin{tcolorbox}[colback=grey!6,colframe=grey!45,boxrule=0.7pt,arc=2pt]
\textbf{Nothing was promoted.} Eleven operator contracts and one new \emph{implementation} of an
existing one live in \code{eye\_ops.py} and \code{muscle\_ops.py}; importing them registers them,
and the run is executed by the stock \code{plexus.schema.load} and \code{plexus.engine.run}.
\code{src/plexus} is untouched. The deliverable is a \code{spec.yaml}.
\end{tcolorbox}

\subsection{The decomposition into sets}

\begin{center}\small
\begin{tabular}{@{}llp{7.2cm}@{}}
\toprule
\textbf{set} & \textbf{what it is} & \textbf{state} \\
\midrule
\code{eye} & the globe, as one organ & \code{pos} (centroid), \code{gaze} = (horizontal, vertical, torsion), degrees \\
\code{mpm\_particle} & the globe's material points (parent \code{eye}) & \code{pos}/\code{vel} + $F$, $C$, mass, Lamé, volume \\
\code{muscle} & the six extraocular muscles & \code{act} (the only integrated block), \code{tension}, \code{length} \\
\code{muscle\_particle} & each muscle's material points (parent \code{muscle}) & as \code{mpm\_particle} \\
\code{orbit} & the bony socket & \code{pos}, the centre of the cup \\
\code{mpm\_grid} (field) & \textbf{one} background grid, shared by both bodies & --- \\
\bottomrule
\end{tabular}
\end{center}

The shared grid is the design. It is what mechanically couples a contracting muscle to the sclera,
and it is free --- sharing a background grid is simply how MLS-MPM couples any two bodies. The
stock \code{mpm\_scatter} \emph{overwrites} the grid, so a second \emph{implementation} of that
same contract, \code{implementation: accumulate}, was registered: identical biology, different
numerics. Per Plexus\,2 that is the correct way to extend --- never an edit to the engine.

\subsection{The operators}

\begin{center}\small
\begin{tabular}{@{}llp{8.0cm}@{}}
\toprule
\textbf{operator} & \textbf{kind} & \textbf{what it does} \\
\midrule
\code{eye\_anatomy} & rewire & squash the seeded ball into the ovoid; label sclera / cornea / iris / pupil / choroid / vitreous / lens; per-region Lamé \\
\code{muscle\_morphogenesis} & rewire & shape each muscle's points into a tapered strap along its arc of contact; fibres; bone cap and tendon cap \\
\code{eye\_pose} & aggregate & Kabsch fit of the shell against its rest state $\rightarrow$ gaze and torsion \\
\code{muscle\_geometry} & aggregate & bin points by fibre coordinate $\rightarrow$ centreline, \textbf{length}, insertion, line of action, rotation axis \\
\code{oculomotor\_drive} & lateral & gaze program $\rightarrow$ PID $\rightarrow$ rectified projection onto each muscle's axis $\rightarrow$ activation dynamics \\
\code{muscle\_contract} & exchange & $\sigma_a = A\,a\,(1+\beta(\lambda-1))\,\mathbf{f}\mathbf{f}^{\!\top}$ along the fibre \\
\code{bone\_anchor} & lateral & pin the origin cap to the skull \\
\code{muscle\_sleeve} & lateral & the orbital pulley: a purely \emph{transverse} restoring force, so the path is held but shortening is free \\
\code{orbit\_socket} & lateral & penalty contact with the spherical cup + orbital-fat suspension \\
\code{mpm\_scatter}\,[\code{accumulate}] & exchange & the second body into the shared grid \\
\bottomrule
\end{tabular}
\end{center}

The schedule closes a loop: pose readout $\rightarrow$ gaze error $\rightarrow$ innervation
$\rightarrow$ contraction $\rightarrow$ MLS-MPM $\rightarrow$ motion $\rightarrow$ pose readout.

%======================================================================
\section{Where we actually are}

\begin{center}\small
\begin{tabular}{@{}llp{7.6cm}@{}}
\toprule
\textbf{phase} & \textbf{state} & \textbf{what it produced} \\
\midrule
0 --- build the plant & \dn & 25 trials, 4 named ceilings, a working point, a metric suite \\
1a --- differentiable identification & \wrong & a surrogate that fails its own fidelity test, and refuses to report \\
1b --- the dynamics-watching critic & \wrong & a VLM that reads a rotating eye as static \\
2 --- PID autotune by step injection & \dn & the $3\times6$ static gain matrix: \textbf{mechanics-limited by $7\times$}, and the obliques named as the binding constraint \\
3 --- buy the authority back, one factor at a time & \fl & what each Phase-0 fix is actually worth, in degrees per activation \\
\bottomrule
\end{tabular}
\end{center}

Two automatic instruments were built to judge the plant and \textbf{both failed}. That is why
Phase~2 is an \emph{open-loop measurement} rather than another optimisation --- and why it is the
first instrument in this campaign that passes a fidelity check instead of failing one.

\begin{tcolorbox}[colback=warn!6,colframe=warn!50,boxrule=0.7pt,arc=2pt]
\textbf{The archive was pruned to the Phase-2 baseline.} Only \code{t01}, \code{t02} and
\code{t03\_c\_a} survive on disk; the run folders from \code{t04} onward were deleted after the
campaign closed. Their metrics are in the git history and are quoted here as the record.
\end{tcolorbox}

%======================================================================
\section{Phase 0 --- Build the plant}

\begin{phasebox}{Phase 0 --- define \hfill \dn}
\textbf{Goal.} A globe that deforms a little but does not come apart, six muscles that visibly
contract, a skull that keeps the eye in its socket, and a movie showing deformation, stress and
activation --- all of it a \code{spec.yaml}, with the operators registered from the prototype
folder and the engine untouched.

\textbf{The constraint that changed the model.} The first version had a muscle as a \emph{line of
action}, its pull a body force injected at the insertion patch. Nothing contracted, nothing
deformed, and the tendon's grip on the sclera was \emph{asserted}. Making the muscles tissue ---
their own particle set, coupled through the shared grid --- is what turned a picture into a
mechanism.

\textbf{What ``done'' means, written as a test.} \code{run\_eye.verdict}: finite; reaches its
commands within $6^\circ$; torsion $>6^\circ$; horizontal range $>35^\circ$; recruitment 4/4;
centroid drift $<6\%$ of the globe radius; globe strain$_{99}$ in $[0.004, 0.12]$; peak shortening
in $[4\%, 38\%]$; globe radius worst $<6\%$ and roundness spread $<9\%$.
\end{phasebox}

\begin{runbox}{Phase 0 --- run \hfill 25 trials, 4 ceilings}
Each ceiling was a real mechanism, not a tuning miss. They are worth more than the working point.

\footnotesize
\begin{tabular}{@{}p{3.0cm} p{11.0cm}@{}}
\toprule
\textbf{ceiling} & \textbf{what it turned out to be}\\
\midrule
$A$ and $E$ are not independent (t02) & A muscle shortens until passive tension balances active
stress, so steady shortening is set by $A/E$ \emph{alone} while delivered force is
$A\times$cross-section. Raising $A$ by itself: \textbf{88\% shortening, globe crushed.}\\
$A/E$ is a cliff, not a trade-off (t03--t04) & 0.25 clean, 0.30 collapses --- a \emph{buckling}
threshold. But the stiff passive element that bought stability made the antagonist so stiff it ate
most of the agonist's force: \textbf{the eye capped at $11^\circ$.}\\
Length buys gaze (t06--t08) & $\theta_{\max}\!\sim\!(A/E)(L/R)$. Truncating each muscle to 55\% of
its path --- a convenience taken so the four recti would not project onto one blob --- had left the
recti at $L/R=2.3$ and the obliques at 1.15, against $\approx3.3$ in a real orbit. Restoring the
length and spreading the origins around the annulus of Zinn (a \emph{ring}, not a point):
\textbf{$12.7^\circ \to 29.6^\circ$.}\\
The missing mechanism was the pulley (t15--t24) & A free-floating strap is a column in compression
the moment it contracts. \code{muscle\_sleeve} restores only the component \emph{perpendicular} to
the local fibre, tapering to zero over the distal third --- Demer's active pulley, which is exactly
a transverse constraint. Everything else held fixed, it raises the buckling ceiling from
$\approx$0.25 to \textbf{between 0.28 (clean) and 0.32 (folds)}.\\
\midrule
\emph{and a material error} (t18) & The vitreous at $E=9$ is a shear modulus of 3.8 --- very nearly
a fluid. The fixed-corotated law has no failure criterion, so under a hard tendon pull the interior
flowed and \textbf{the globe came apart}. Interior raised to 45 / 130 against a 420 sclera: same
$\approx9\times$ core-to-shell contrast, enough shear modulus \emph{everywhere} to hold together.\\
\bottomrule
\end{tabular}
\normalsize

\textbf{Working point} (t23/t25, the full 1830-frame programme): range $33.5/18.3/7.5^\circ$,
shortening 24--31\% with no buckling, centroid drift 1.8\% of the globe radius,
strain$_{99}$ 0.112, and the globe's own radius mean $-0.04\%$, worst $0.22\%$, roundness spread
$0.86\%$. \emph{The eye deforms locally and holds its shape globally} --- which is the whole reason
to use MPM here, and is now a number rather than a look at the movie.

\textbf{Six of nine acceptance checks pass.} \code{reaches\_commands} fails (max settled error
$17.1^\circ$); \code{wide\_gaze\_range} fails narrowly ($33.5^\circ$ against a $35^\circ$ bar); and
\code{correct\_recruitment} fails \emph{on the atlas programme only}, at 14/26 --- that programme's
26 holds include diagonal combinations for which no single muscle is ``the'' expected agonist, so a
top-1 test is the wrong test there. On the four-hold probe programme it is 4/4 in every
mechanically-stable trial. The real gap is the first one: the eye settles at roughly 60\% of a
$25^\circ$ command, and everything from here is about that.
\end{runbox}

\textbf{A lesson about pictures.} The muscle truncation was made so the anatomy would \emph{render}
well. It cost $17^\circ$ of gaze. A convenience taken for the figure was a mechanical error, and
nothing in the pipeline would have caught it --- the movie looked better after the change.

%======================================================================
\section{Phase 1 --- Can an automatic instrument judge this?}

Two were built. Both failed, and both failures are kept because they are the reason Phase~2 is
shaped the way it is.

\subsection{Phase 1a --- differentiable identification and gain tuning}

\begin{phasebox}{Phase 1a --- define \hfill \dn}
\textbf{Claim to test.} The gap between commanded and achieved gaze is a \emph{controller} problem,
and gradient descent on the gains will close it.

\textbf{Method, and why not the obvious one.} The engine can keep the tape
(\code{run(\ldots, grad=True)}), but a rollout of 45k$+$13k material points $\times$ 25 substeps
$\times$ hundreds of frames unrolls into a graph whose memory scales as points~$\times$~steps ---
exactly the cost Plexus\,2 says makes a global loss impractical. \code{../inverse\_slime/} gets
around it by refusing to roll out at all, so \code{tune\_gaze.py} does the same in two stages:
identify $\ddot\theta = Bu - C\dot\theta - K\theta$ per axis by ridge least squares from an archived
run, then run Adam on $(k_p,k_i,k_d,\text{gain},\text{tonic})$ over that differentiable surrogate,
and verify back in the full simulation.

\textbf{Success looks like} a surrogate that reproduces the run it was fitted to, and a tuned gain
set that measurably improves tracking in the real simulation.
\end{phasebox}

\begin{runbox}{Phase 1a --- run \hfill the instrument failed its own check}
\textbf{First the derivatives, not the model.} $R^2$ on $\ddot\theta$ came out 0.28--0.40.
Differentiating a strided, noisy signal \emph{twice} amplifies its noise by $\sim h^{-2}$; reading
the derivatives off a local Savitzky--Golay polynomial instead raised it to \textbf{0.48--0.72} with
no change to the model at all.

\textbf{Then the model.} Tuning still bought $1.008\times$. The first explanation offered --- ``the
drive must be saturated'' --- was \emph{wrong}, and the tool's own diagnostic said so: the agonist
sits above 0.98 on only 6\% of frames. So the surrogate was asked to reproduce the excursion the
real eye actually made, under the gains it actually used:

\begin{center}\small
\begin{tabular}{@{}lccc@{}}
\toprule
peak excursion & horizontal & vertical & torsion \\
\midrule
real            & $17.0^\circ$ & $9.0^\circ$ & $2.8^\circ$ \\
surrogate       & $\ \ 1.7^\circ$ & $3.0^\circ$ & $2.2^\circ$ \\
\bottomrule
\end{tabular}
\end{center}

\begin{tcolorbox}[colback=bad!4,colframe=bad!45,boxrule=0.7pt,arc=2pt]
\textbf{The instrument failed its own check, so the measurement was thrown away.} A good $R^2$ on
$\ddot\theta$ certifies the \emph{transients}; the standing error we wanted to tune away lives in
the steady state $\theta_{ss}=Bu/K$, and the fit misses the real excursion by a factor of ten in the
horizontal. \code{tune\_gaze} now runs this fidelity test and \textbf{refuses to report a tuned
number when it fails} --- the discipline Plexus\,2 puts first: verify the instrument before trusting
the measurement.
\end{tcolorbox}

\textbf{Two defects in the method itself, admitted here because Phase 2 exists to fix them.}

\begin{enumerate}
\item \textbf{Closed-loop identification bias.} The fit was taken from a run \emph{with the
controller closed}. In that data the input $u$ is a function of the state --- that is what feedback
means --- so the regressor is correlated with the residual and least squares is biased. This is a
textbook error and is very likely why $B/K$ predicted $1.7^\circ$ where the eye did $17^\circ$.
\item \textbf{The saturation statistic was the wrong statistic.} ``6\% of frames'' was measured over
\emph{all} frames, most of which are transitions or primary-position holds where the error is near
zero. The number that matters is saturation \emph{of the agonist, during the settled tail of a
hold}; with $k_p\,e = 0.11\times10^\circ$, times gain 1.9, clipped at 1.0, it is very likely pinned
there. \textbf{So the conclusion ``the residual is mechanical, not control'' was not established.}
It rested on a broken surrogate and a badly-chosen statistic, and this note previously stated it
with more confidence than the evidence carried.
\end{enumerate}
\end{runbox}

\subsection{Phase 1b --- the dynamics-watching critic}

\begin{phasebox}{Phase 1b --- define \hfill \dn}
Plexus\,2 names a \emph{dynamics-watching critic} as a required agent, on the grounds that a worker
agent reads a static strip but cannot watch a movie. That is cheap to test here, so: show the local
Gemma-4-12B captioner the anterior panel, cropped so that no axis label, legend or title can leak
the answer, and ask what happens. \textbf{Success is the word ``rotating''.}
\end{phasebox}

\begin{runbox}{Phase 1b --- run \hfill it reads a rotating eye as static}
\textbf{Whole 1830-frame run.} It read the anatomy well --- ``a large, light-grey circular ring'',
``four small yellow clusters arranged symmetrically around the inner black core'', ``elongated
appendages in blue, red, purple and green'', even the dashed cup outline --- and concluded:
\emph{``The overall configuration remains static throughout the sequence of keyframes.''} It never
used the word eye.

That could have been a sampling artefact, since the programme returns to primary position between
oscillation blocks and evenly-spaced keyframes may all land on one pose. So:

\textbf{One excursion only} --- a 1.8\,s clip verified from the recorded trace to span
$0 \to +14.8 \to -16.3^\circ$. This time it counted \emph{six} appendages correctly, and it did
report a change --- the wrong one: \emph{``the color of the appendage on the right side shifts from
a dark olive green to a bright yellow \ldots\ the other colored appendages remain relatively stable
in their positions \ldots\ the central black void remains constant.''}

\begin{tcolorbox}[colback=bad!4,colframe=bad!45,boxrule=0.7pt,arc=2pt]
\textbf{The paper's claim, isolated.} The colour shift it caught is the medial rectus brightening
with activation --- a cue encoded \emph{within} each frame. The $31^\circ$ rotation, encoded
\emph{across} frames, it missed entirely, calling the pupil stationary while the pupil swept a third
of the globe. The failure is not competence at seeing; it is competence at seeing \emph{change}. A
keyframe caption pipeline is structurally unable to referee a rotation, and no prompt fixes it.
Captions in \code{archive/vlm/}.
\end{tcolorbox}

Small irony worth keeping: the gold iridophore flecks exist \emph{because} a round pupil on a
rotating sphere has almost no per-frame signature. The VLM saw the flecks and never tracked their
angle.
\end{runbox}

%======================================================================
\section{Phase 2 --- PID autotune by open-loop step injection}

\begin{phasebox}{Phase 2 --- define \hfill \dn}
\textbf{The question, stated so it can come out either way.} Is the standing gaze error a
\emph{control} problem or a \emph{mechanical} one? Phase~1a asserted an answer with a broken
instrument. This phase measures it.

\textbf{Method --- the standard autotune, which is what Phase 1a should have been.} Open the loop
and inject a step on \emph{one muscle at a time}:

\begin{itemize}
\item a new operator \code{muscle\_probe} writes activation from a prescribed waveform and ignores
the gaze error entirely, so the input is no longer a function of the state and the
closed-loop bias of Phase~1a is gone by construction;
\item six runs, one per muscle: settle at tonic, step to $a_{hi}$, hold, release;
\item each yields a 3-vector gaze response, so the six together give the full $3\times6$
\textbf{static gain matrix} $G_{km} = \Delta\theta_{k,\infty}/\Delta a_m$ read straight off the
plateau --- the quantity the acceleration fit had no way to see --- plus the rise time and
overshoot per channel.
\end{itemize}

\textbf{Baseline: \code{archive/t03\_c\_a}}, loaded from its own \code{spec.yaml} rather than
rebuilt, so the probe measures exactly the archived configuration. It is the right baseline because
it is \emph{stable}: 4/4 recruitment, strain$_{99}$ 0.052, 17--27\% shortening, no buckling, no
sleeve, $A=60$, $E=240$, and a modest $12.7^\circ$ range. A clean step response wants a
well-behaved plant, not a strong one.

\begin{tcolorbox}[colback=rule!5,colframe=rule!45,boxrule=0.7pt,arc=2pt]
\textbf{The prediction, written down before the run.} Extrapolated to full activation, the lateral
rectus static gain $\Delta\theta_{h,\infty}$ on \code{t03\_c\_a} will land in
$\mathbf{8^\circ}$--$\mathbf{16^\circ}$ --- \emph{below} the $26^\circ$ it is being commanded to
reach. If that holds, this configuration is mechanically incapable of the command and no PID retune
can close the gap. \\[3pt]
\textbf{The falsifier.} If $\Delta\theta_{h,\infty} \ge 26^\circ$, the plant has the authority and
the loop is throwing it away --- the gap is control, Phase~1a's conclusion was wrong, and the
autotune should recover most of it. \\[3pt]
\textbf{Secondary prediction.} LR and MR will show a nearly pure direct-axis gain (off-axis
$<15\%$), while SR, IR, SO and IO will show substantial cross-axis terms --- the same
primary/secondary/tertiary structure the geometry already predicts, now measured in
degrees-per-activation instead of as an axis direction.
\end{tcolorbox}

\textbf{Then, and only then, tune.} The identified model is a handful of numbers, so the gains come
out in milliseconds --- no gradient through MPM needed. Two things to get right when it happens:
the actuator saturates and the drive is rectified, so this is a \emph{Hammerstein} system (static
nonlinearity then linear dynamics) and the nonlinearity must be fitted explicitly rather than
absorbed into the linear part; and the integrator needs anti-windup, which \code{i\_clip} provides
but which has never been exercised.

\textbf{Where differentiability actually earns its keep.} Not here. The reduced plant is fitted, not
differentiated. The tape is worth its cost for tuning the \emph{anatomy} --- insertion angles, arc
of contact, pulley placement --- where there is no reduced model to fit. That is the more
interesting inverse problem and it is untouched.
\end{phasebox}

\begin{runbox}{Phase 2 --- run \hfill \dn\ (6 of 6, archived \code{t04}--\code{t09})}
Six open-loop steps from tonic to full activation, 650 frames each, held 500. The plateau is
reached by frame $\approx$400 after a 40\% overshoot and a ring-down, so the settled tail is a real
static gain and not a transient sample --- which every earlier measurement in this campaign had
been.

\textbf{The static gain matrix} (degrees of gaze per unit activation, all other muscles at tonic):

\begin{center}\small
\begin{tabular}{@{}lrrrrrr@{}}
\toprule
 & \textbf{LR} & \textbf{SR} & \textbf{MR} & \textbf{IR} & \textbf{SO} & \textbf{IO}\\
\midrule
horizontal & $4.25$ & $-2.18$ & $\mathbf{-11.65}$ & $-2.51$ & $-0.20$ & $-0.22$\\
vertical   & $-0.04$ & $4.87$ & $0.04$ & $-4.85$ & $-0.78$ & $0.52$\\
torsion    & $-0.08$ & $2.06$ & $-0.07$ & $-2.00$ & $\mathbf{1.10}$ & $-0.59$\\
\midrule
off-axis fraction & 0.02 & 0.62 & 0.01 & 0.66 & 0.73 & 0.97\\
$t_{63}$ (s) & 0.180 & 0.135 & 0.165 & 0.135 & 0.210 & 0.225\\
overshoot & 42\% & 45\% & 35\% & 45\% & 30\% & 35\%\\
\bottomrule
\end{tabular}
\end{center}

\begin{tcolorbox}[colback=bad!4,colframe=bad!45,boxrule=0.7pt,arc=2pt]
\textbf{The prediction is broken, and the conclusion is stronger for it.} The registered band was
$8$--$16^\circ$; the lateral rectus at full activation reaches $\mathbf{3.7^\circ}$ against a
$26^\circ$ command. Not a marginal miss --- \textbf{mechanics-limited by a factor of seven}. No
value of any gain closes that, and the Phase~1a conclusion, which happened to say the same thing,
was right for none of the reasons it gave.
\end{tcolorbox}

\textbf{The secondary prediction held, and sharply.} LR and MR come out essentially pure
direct-axis (off-axis 0.02 and 0.01 against a predicted $<0.15$), while SR, IR, SO and IO are all
substantially coupled (0.62--0.97). SR and IR are near-perfect mirror images, and SR's
torsion-to-elevation ratio of $2.06/4.87 = 0.42$ matches the ratio $0.39/0.90 = 0.43$ that the
\emph{geometry} predicted in \S6, from a completely independent calculation.

\textbf{Two things the measurement found that nobody predicted.}

\begin{enumerate}
\item \textbf{MR is 2.7$\times$ stronger than LR} ($-11.65$ against $+4.25$), though their peak
tensions differ by 5\%. The medially-tilted apex makes the medial rectus' pull far more tangential
than the lateral rectus', and the arc-of-contact wrap does the rest. This was confirmed
independently, from data collected before the probe existed: \code{t03\_c\_a}'s closed-loop trace
runs $+2.62^\circ$ to $-10.05^\circ$ --- the eye adducts four times further than it abducts, and
nobody had noticed.
\item \textbf{The obliques are the weakest muscles in the plant, and the vertical recti are better
torters than they are.} SO produces $1.10^\circ$ of torsion per unit activation against SR's
$2.06^\circ$. That is anatomically wrong, and it traces straight back to the Phase-0 length
ceiling: at \code{frac}=0.55 the obliques' post-pulley path is 0.132 against the lateral rectus'
0.242, so $L/R\approx1.15$. \emph{This is why torsion was so hard to command all through Phase 0},
and it is the first quantitative statement of it.
\end{enumerate}

\begin{tcolorbox}[colback=good!5,colframe=good!45,boxrule=0.7pt,arc=2pt]
\textbf{The instrument passes the check that Phase 1a failed.} Assume the closed loop saturates the
agonist and relaxes the antagonist off tonic; the open-loop matrix then \emph{predicts} a
closed-loop adduction of $-10.61^\circ$, against $-10.05^\circ$ actually observed in
\code{t03\_c\_a} --- \textbf{5\%}, on a run the matrix was never fitted to. (The other excursions
come in below prediction, as they must: full saturation is an upper bound and the PID rarely
reaches it.) Phase~1a's surrogate predicted $1.7^\circ$ where the eye did $17.0^\circ$. That is the
difference between an open-loop and a closed-loop identification, and it is the whole reason this
phase exists.
\end{tcolorbox}

\textbf{So the PID autotune is answered before it is run:} there is nothing for it to recover on
this configuration. The dynamics it would need are now measured anyway ($t_{63}\approx0.14$--$0.23$\,s,
30--45\% overshoot, so $\zeta\approx0.3$), and they are worth keeping for when the plant can
actually reach its commands.
\end{runbox}

%======================================================================
\section{Phase 3 --- Buy the authority back, one factor at a time}

\begin{phasebox}{Phase 3 --- define \hfill \fl}
\textbf{Not a controller.} Phase~2 settled that: this plant is short of authority by $7\times$, and
no gain recovers what the mechanics does not produce. So Phase~3 goes after the mechanics --- and
uses the Phase-2 probe as its instrument, because that probe is the first thing in this campaign
that passed a fidelity check.

\textbf{The three interventions, and why these three.} Phase~0 found them by closed-loop guessing
and never learned what any one of them was worth. \code{t03\_c\_a} predates all three, so it is the
clean starting point:

\begin{enumerate}
\item \textbf{Length} --- \code{frac} $0.55\to0.95$ and the sclera stand-off $0.020\to0.042$.
$\theta_{\max}\sim(A/E)(L/R)$, and this is the term the baseline is worst on. Measured: $L$ goes
$0.242\to0.449$ for LR and $0.132\to0.253$ for SO, so $L/R$ does roughly double.
\item \textbf{The pulley} --- add \code{muscle\_sleeve}. An anti-buckling device, not a force
multiplier.
\item \textbf{Drive} --- $A/E$ $0.25\to0.28$, the top of the range the pulley was measured to allow
in Phase~0.
\end{enumerate}

\begin{tcolorbox}[colback=bad!4,colframe=bad!45,boxrule=0.7pt,arc=2pt]
\textbf{A fourth intervention was dropped, because it turned out to be already applied --- and
finding that out exposed a reproducibility defect.} Spreading the rectus origins around the annulus
of Zinn was on this list until the baseline was checked. \code{muscle\_morphogenesis} reads the
origins from \code{eye\_anatomy.origins\_world()} at \emph{run} time, so
\textbf{an archived \code{spec.yaml} does not fully determine its own run}: \code{t03\_c\_a} was
generated before \code{ANNULUS\_RING} existed and recorded a rest length of $0.259$ for the lateral
rectus, but re-running that same file today gives $0.242$. The Phase-2 baseline therefore already
carries the spread origins, the intervention's effect was never measured, and it cannot be recovered
without reverting the module. \emph{Geometry a spec does not carry is geometry the spec cannot
reproduce} --- and every Phase-2 number in this note is against the re-run baseline, not against the
original \code{t03\_c\_a}.
\end{tcolorbox}

\textbf{One at a time, re-probing after each.} Six probe runs per intervention, 18 in total.
That is the discipline Phase~0 broke when it changed three things at once and then could not say
which had caused a divergence.

\begin{tcolorbox}[colback=rule!5,colframe=rule!45,boxrule=0.7pt,arc=2pt]
\textbf{The predictions, written down before the runs.}
\begin{itemize}
\item \textbf{Length.} $L/R$ roughly doubles for every muscle, so if $\theta_{\max}\sim L/R$ holds,
the gains roughly double: LR horizontal $4.25 \to \mathbf{7}$--$\mathbf{9}$, SO torsion
$1.10 \to \mathbf{1.9}$--$\mathbf{2.4}$ deg per unit activation.
\item \textbf{The pulley, the sharp one.} If \code{muscle\_sleeve} is what this note claims it is
--- a purely transverse constraint that prevents buckling and nothing else --- it must leave the
\emph{static} gain matrix nearly unchanged: \textbf{under 20\% on every entry}. If instead the gains
move a lot, the sleeve is doing something we have not described, and the Phase-0 story about it is
wrong.
\item \textbf{Drive.} Gains are linear in $A$, so $+12\%$, no more.
\item \textbf{The headline, and the uncomfortable one.} Compounded, LR reaches
$3.7 \times 2 \times 1.12 \approx \mathbf{8}$--$\mathbf{9^\circ}$ --- \textbf{still nowhere near the
$26^\circ$ it is commanded to make}. Prediction: \emph{all three Phase-0 fixes together do not make
this plant able to reach its commands}, and what has to change is either the command amplitude or
something more basic than a parameter --- cross-section, globe compliance, or how the tendon
couples.
\end{itemize}
\end{tcolorbox}

\textbf{Why the last prediction is worth making even though it is unwelcome.} Phase~0 reported
$29.6^\circ$ of closed-loop range after the length fix and read that as success. Phase~2 has since
shown that a closed-loop range is roughly \emph{twice} the single-muscle static gain, because the
antagonist relaxes as the agonist pulls --- so $29.6^\circ$ of range is about $15^\circ$ of
authority per direction, against commands of $26^\circ$. The gap was always there. It was hidden by
measuring the wrong quantity.
\end{phasebox}

\begin{runbox}{Phase 3 --- run \hfill \todo}
Three interventions, six probe runs each, and the gain matrix after every one. This box is where
the answer goes --- including, and especially, if the headline prediction above is wrong and the
plant does reach its commands.
\end{runbox}

%======================================================================
\section{The one result worth defending}

\textbf{No muscle action is tabulated anywhere in this prototype.} Each muscle's rotation axis is
re-measured every frame from where its tissue actually is, as $\hat{\mathbf n}\times\hat{\mathbf u}$
(insertion radius $\times$ line of action), and \code{oculomotor\_drive} simply rectifies the
projection of the desired rotation onto it --- Sherrington's law of reciprocal innervation, one line
of code. So the textbook primary/secondary/tertiary actions have to \emph{emerge} from where each
muscle inserts and where it pulls from. They do:

\begin{center}\small
\begin{tabular}{@{}llll@{}}
\toprule
\textbf{muscle} & \textbf{measured axis $(x,y,z)$} & \textbf{reading} & \textbf{textbook} \\
\midrule
LR & $(\,0.03,\ \ 1.00,\ -0.05)$ & pure abduction & abduction only \\
MR & $(\,0.02,\ -1.00,\ \ 0.03)$ & pure adduction & adduction only \\
SR & $(-0.90,\ -0.17,\ \ 0.39)$ & elevation $>$ intorsion $>$ adduction & same order \\
IR & $(\,0.90,\ -0.22,\ -0.37)$ & depression $>$ extorsion $>$ adduction & same order \\
SO & $(\,0.61,\ \ 0.14,\ \ 0.78)$ & intorsion $>$ depression $>$ abduction & same order \\
IO & $(-0.66,\ \ 0.18,\ -0.73)$ & extorsion $>$ elevation $>$ abduction & same order \\
\bottomrule
\end{tabular}
\end{center}

All six agree, in sign and in rank. The medial tilt of the orbital apex is what gives the vertical
recti their torsional and horizontal components; a co-axial apex makes SR and IR pure elevators and
depressors, which is anatomically wrong. Recruitment measured from the runs agrees with this table
on 4 of 4 commands in every mechanically-stable trial of the probe programme. \emph{Not} on the
atlas programme (14/26) --- but that is a defect of the test, not of the plant: half its holds are
diagonal commands recruiting two muscles at once, against which a top-1 comparison is meaningless.

\begin{tcolorbox}[colback=good!5,colframe=good!45,boxrule=0.7pt,arc=2pt]
\textbf{Why this is worth more than a nice movie.} It is a falsifiable prediction that geometry
alone determines the actions, tested against a table nobody consulted while writing the code. It is
also the Plexus point in miniature: the mechanism is in the operators and the anatomy, not in a
lookup table inside a simulator. Phase~2 sharpens it --- the same claim, in degrees per unit
activation.
\end{tcolorbox}

%======================================================================
\section{Two modelling decisions that are not obvious}

\textbf{The globe is an ovoid; the cup is a sphere.} Teleost eyes are flattened along the optic
axis. The \code{mpm\_particle} entity can only seed a uniform ball, so the rest configuration is
that ball affinely squashed, $z\to kz$ with $k=0.82$. An affine image of a uniform ball is a uniform
ellipsoid, so density stays uniform, and because this is an \emph{initial condition} and not a
deformation, $F$ stays the identity --- the ovoid is undeformed and carries no residual stress. The
bony cup stays spherical with a radius just above the equatorial semi-axis: an ovoid cannot rotate
inside a socket that matches it, and in the animal the gap is filled by orbital fat, not bone.

\textbf{The fat exerts no torque.} The retrobulbar fat pad is a restoring force toward the socket
centre applied \emph{uniformly} to every point of the globe. A uniform body force produces no torque
about the centroid, so the fat recentres the eye without ever resisting gaze --- exactly the division
of labour the anatomy has. Without it, the residual normal component of muscle tension slowly walks
the globe around inside the cup.

\textbf{Both ends of every muscle are fixed, by two different mechanisms.} The origin cap is pinned
explicitly by \code{bone\_anchor}, a penalty force toward its frame-0 rest positions. The insertion
is pinned by nothing: the tendon cap is given a \emph{negative} stand-off so it overlaps the sclera,
and points that share a B-spline stencil node share one grid velocity --- a no-slip weld, for free.
The honest caveat is that this weld is resolution-dependent and \emph{cannot slip}, whereas a real
tendon slides in its sheath. That is why the bellies are held about five cells clear of the globe;
otherwise the whole arc of contact would weld and freeze the rotation.

%======================================================================
\section{The movie}

Eight panels, black ground, lower-case parenthesised labels, one shared time cursor.

\begin{center}\small
\begin{tabular}{@{}llp{8.4cm}@{}}
\toprule
a) & anterior view & the cosmetic eye --- white sclera, silvery iris, large black pupil, gold iridophore flecks, after the reference photograph --- with the six muscles as their own material points, brightening with activation \\
b) & lateral view & the ovoid profile in the bony cup, the obliques, the trochlea \\
c) & strain & Green--Lagrange $\|E\|$ on a cut through globe \emph{and} muscles \\
d) & stress & von Mises on the same cut \\
e) & grid momentum & $|v|$ on the shared MLS-MPM grid --- the medium through which a contracting muscle reaches the sclera \\
f) & activation & the six innervations against time \\
g) & length & each muscle as \% of rest, with \textbf{the globe's own radius} on the same axis: the muscles shorten, the eye must not \\
h) & gaze & horizontal / vertical / torsion against command, with the tracking error and its running mean, and the per-axis RMS printed \\
\bottomrule
\end{tabular}
\end{center}

Panels c--e follow the codebase's \code{grid.mp4} diagnostic, extended to two coupled bodies. Two
rendering notes that are not cosmetic: the far hemisphere is culled, so the eye reads as a lit solid
rather than a point cloud; and muscle points falling inside the globe's exact orthographic
silhouette \emph{and} behind its centre are hidden, because a depth-sorted scatter alone lets a
muscle show through the gaps between globe particles --- most visibly straight across the black
pupil.

%======================================================================
\section{What is deliberately not done}

\begin{itemize}
\item \textbf{No promotion.} Eleven contracts and one implementation sit in \code{prototype/}.
Plexus\,2 requires evidence of reuse beyond the originating prototype first, and only
\code{muscle\_sleeve} and \code{mpm\_scatter}\,[\code{accumulate}] look plausibly general.
\item \textbf{\code{muscle\_traction} and \code{muscle\_insertion} are dead code, kept.} They are
the \emph{first} version --- a line of action with a body force at the insertion. This note argues
against them, and the argument should be checkable.
\item \textbf{The integral term is implemented but off by default.} It is the oculomotor neural
integrator and it is correct biology, but it was introduced in the same run as two other changes and
contributed to a divergence. It has not been isolated since, and Phase~2 is where it should be.
\item \textbf{No inverse pass on the anatomy.} Recovering the six insertion geometries from a gaze
trace is the interesting Loop~II question here, and it is untouched.
\end{itemize}

\end{document}
